\begin{center}\textit{Einige spezielle Sätze.}\end{center}
\SixthArticle{53}
Alle bisher von uns hergeleiteten Beziehungen sind insofern von größter Allgemeinheit, als die Elemente $\alpha,\beta,\gamma$ keinen Bedingungen unterliegen. Darüber hinaus haben wir jedoch mehrere andere gefunden, die besondere Bedingungen zwischen den Elementen $\alpha,\beta,\gamma$ voraussetzen: Zweifellos liegen noch weit mehr verborgen, und vielleicht lassen sich künftig selbst die hier mitgeteilten aus höheren Prinzipien herleiten.
\par\SourceLeaf{233}
Setzen wir zunächst in Gleichung 80 $x=\frac{4y}{(1+y)^2}$, so ist
\[
\mathrm dx=\mathrm dy\mathbin{.}\frac{4(1-y)}{(1+y)^3}
\]
und somit
\[
\begin{aligned}
\frac{\mathrm dP}{\mathrm dx}&=\frac{\mathrm dP}{\mathrm dy}\mathbin{.}\frac{(1+y)^3}{4(1-y)}\\
\frac{\mathrm{dd}P}{\mathrm dx}&=\mathrm d\frac{\mathrm dP}{\mathrm dy}\mathbin{.}\frac{(1+y)^3}{4(1-y)}\\&\qquad+\frac{\mathrm dP}{\mathrm dy}\mathbin{.}\frac{(2-y)(1+y)^2}{2(1-y)^2}\mathrm dy\\
\frac{\mathrm{dd}P}{\mathrm dx^2}&=\frac{\mathrm{dd}P}{\mathrm dy^2}\mathbin{.}\frac{(1+y)^6}{16(1-y)^2}\\&\qquad+\frac{(2-y)(1+y)^5}{8(1-y)^3}\mathbin{.}\frac{\mathrm dP}{\mathrm dy}
\end{aligned}
\]
Dadurch wird jene Gleichung zu
\[
\begin{aligned}
0={}&\alpha\beta P\\
&-(\gamma(1+y)^2-4(\alpha+\beta+1)y)\frac{1+y}{4(1-y)}\mathbin{.}\frac{\mathrm dP}{\mathrm dy}\\
&-\frac{\mathrm{dd}P}{\mathrm dy^2}\mathbin{.}\frac{y(1+y)^2}4
-\frac{y(2-y)(1+y)}{2(1-y)}\mathbin{.}\frac{\mathrm dP}{\mathrm dy}
\end{aligned}
\]
oder
\[
\begin{aligned}
&0=4\alpha\beta(1-y)P\\
&-(\gamma(1+y)^2-4(\alpha+\beta+1)y-2y(2-y)(1+y))\frac{\mathrm dP}{\mathrm dy}\\
&-(y-yy)(1+y)^2\frac{\mathrm{dd}P}{\mathrm dy^2}\\
&0=4\alpha\beta(1-y)P\\
&-(1+y)(\gamma-(4\alpha+4\beta-2\gamma)y+(\gamma-2)yy)\frac{\mathrm dP}{\mathrm dy}\\
&-(1+y)^2(y-yy)\frac{\mathrm{dd}P}{\mathrm dy^2}
\end{aligned}
\]
Setzt man $P=(1+y)^{2\alpha}Q$, so folgt daraus
\[
\begin{aligned}
\mathrm I.\quad 0={}&2\alpha(2\beta-\gamma+(2\alpha+1-\gamma)y)Q\\
&-(\gamma-(4\beta-2\gamma)y+(\gamma-4\alpha-2)yy)\frac{\mathrm dQ}{\mathrm dy}\\
&-(y-yy)(1+y)\frac{\mathrm{dd}Q}{\mathrm dy^2}
\end{aligned}
\]
Nimmt man nun $\beta=\alpha+\frac12$ an, so erhält diese Gleichung die folgende Form:
\[
\begin{aligned}
0={}&2\alpha(2\alpha+1-\gamma)Q\\
&-(\gamma-(4\alpha+2-\gamma)y)\frac{\mathrm dQ}{\mathrm dy}\\
&-(y-yy)\frac{\mathrm{dd}Q}{\mathrm dy^2}
\end{aligned}
\]
Deren Lösung ist
\[
Q=F(2\alpha,2\alpha+1-\gamma,\gamma,y)
\]
\SourceLeaf{234}
sodass sich ergibt:
\[
\begin{aligned}&[100]\quad(1+y)^{2\alpha}F(2\alpha,2\alpha+1-\gamma,\gamma,y)\\&=F(\alpha,\alpha+\tfrac12,\gamma,\tfrac{4y}{(1+y)^2})\end{aligned}
\]
\SixthArticle{54}
Nehmen wir statt der Beziehung $\beta=\alpha+\frac12$ die Beziehung $\gamma=2\beta$ an, so wird Gleichung I des vorigen Artikels zu
\[
\begin{aligned}
0={}&2\alpha(2\alpha+1-2\beta)yQ\\
&-(2\beta-(4\alpha+2-2\beta)yy)\frac{\mathrm dQ}{\mathrm dy}\\
&-y(1-yy)\frac{\mathrm{dd}Q}{\mathrm dy^2}
\end{aligned}
\]
Setzt man nun $yy=z$, so ist
\[
\begin{aligned}
\frac{\mathrm dQ}{\mathrm dy}&=2y\frac{\mathrm dQ}{\mathrm dz}\\
\frac{\mathrm{dd}Q}{\mathrm dy^2}&=4yy\frac{\mathrm{dd}Q}{\mathrm dz^2}+\frac{2\mathrm dQ}{\mathrm dz}
\end{aligned}
\]
und somit
\[
\begin{aligned}
0={}&\alpha(\alpha+\tfrac12-\beta)Q\\
&-(\beta+\tfrac12-(2\alpha+\tfrac32-\beta)z)\frac{\mathrm dQ}{\mathrm dz}\\
&-(z-zz)\frac{\mathrm{dd}Q}{\mathrm dz^2}
\end{aligned}
\]
Deren Lösung ist
\[
Q=F(\alpha,\alpha+\tfrac12-\beta,\beta+\tfrac12,z)
\]
und daher haben wir
\[
\begin{aligned}&[101]\quad(1+y)^{2\alpha}F(\alpha,\alpha+\tfrac12-\beta,\beta+\tfrac12,yy)\\&=F(\alpha,\beta,2\beta,\tfrac{4y}{(1+y)^2})\end{aligned}
\]
\SixthArticle{55}
Setzen wir zweitens $x=4y-4yy$, so ist
\[
\begin{aligned}
\mathrm dx&=4\mathrm dy(1-2y)\\
\frac{\mathrm dP}{\mathrm dx}&=\frac{\mathrm dP}{\mathrm dy}\mathbin{.}\frac1{4(1-2y)}\\
\frac{\mathrm{dd}P}{\mathrm dx^2}&=\frac{\mathrm{dd}P}{\mathrm dy^2}\mathbin{.}\frac1{16(1-2y)^2}
+\frac{\mathrm dP}{\mathrm dy}\mathbin{.}\frac1{8(1-2y)^3}
\end{aligned}
\]
und Gleichung 80 wird dadurch zu
\SourceLeaf{235}
\[
\begin{aligned}
&0=4\alpha\beta P\\
&-(\gamma-(4\alpha+4\beta+2)y+(4\alpha+4\beta+2)yy)\frac1{(1-2y)}\mathbin{.}\frac{\mathrm dP}{\mathrm dy}\\
&-(y-yy)\frac{\mathrm{dd}P}{\mathrm dy^2}
\end{aligned}
\]
Um den Bruch im zweiten Glied zu beseitigen, muss man $\gamma=\alpha+\beta+\frac12$ setzen; daraus ergibt sich
\[
\begin{aligned}
0={}&4\alpha\beta P\\
&-(\alpha+\beta+\tfrac12-(2\alpha+2\beta+1)y)\frac{\mathrm dP}{\mathrm dy}\\
&-(y-yy)\frac{\mathrm{dd}P}{\mathrm dy^2}
\end{aligned}
\]
Deren Lösung ist
\[
P=F(2\alpha,2\beta,\alpha+\beta+\tfrac12,y)
\]
und daher haben wir
\[
\begin{aligned}&[102]\quad F(\alpha,\beta,\alpha+\beta+\tfrac12,4y-4yy)\\&=F(2\alpha,2\beta,\alpha+\beta+\tfrac12,y)\end{aligned}
\]
Würden wir in dieser Gleichung $y$ durch $1-y$ ersetzen, so ergäbe sich
\[
F(\alpha,\beta,\alpha+\beta+\tfrac12,4y-4yy)=F(2\alpha,2\beta,\alpha+\beta+\tfrac12,1-y)
\]
Daraus scheint das Paradoxon zu folgen:
\[
F(2\alpha,2\beta,\alpha+\beta+\tfrac12,y)=F(2\alpha,2\beta,\alpha+\beta+\tfrac12,1-y)
\]
Diese Gleichung ist gewiss falsch. Um dies aufzulösen, müssen wir uns daran erinnern, sorgfältig zwischen den beiden Bedeutungen des Symbols $F$ zu unterscheiden: Es stellt nämlich \textit{entweder} die Funktion dar, deren Beschaffenheit durch die Differentialgleichung 80 ausgedrückt wird, \textit{oder} allein die Summe der unendlichen Reihe. Die letztere liefert stets eine vollständig bestimmte Größe, solange das vierte Element zwischen $-1$ und $+1$ liegt; Sie müssen jedoch darauf achten, diese Grenzen nicht zu überschreiten, da sonst überhaupt keine Bedeutung mehr übrig bleibt. Die erstere Bedeutung stellt dagegen eine allgemeine Funktion dar, die sich stets nach dem Gesetz der Stetigkeit ändert, wenn sich das vierte Element in stetigem Fluss ändert, gleichgültig, ob Sie ihm reelle oder imaginäre Werte geben, sofern Sie stets die Werte 0 und 1 vermeiden. Daraus ist ersichtlich, dass die Funktion im letzteren Sinn für gleiche Werte des vierten Elements ungleiche Werte annehmen kann, nachdem ein Übergang oder vielmehr eine Rückkehr durch imaginäre Größen erfolgt ist; der von der \textit{Reihe} $F$
\SourceLeaf{236}
dargestellte Wert ist nur einer davon. Es liegt daher überhaupt kein Widerspruch darin, dass \textit{irgendein} Wert der Funktion $F(\alpha,\beta,\alpha+\beta+\frac12,4y-4yy)$ gleich $F(2\alpha,2\beta,\alpha+\beta+\frac12,y)$ ist, während ein \textit{anderer} Wert $=F(2\alpha,2\beta,\alpha+\beta+\frac12,1-y)$ wird; und es wäre ebenso widersinnig, daraus auf die Gleichheit dieser Werte zu schließen, wie wenn man aus $\mathrm{Arc.\ sin}\frac12=30^\circ$, $\mathrm{Arc.\ sin}\frac12=150^\circ$ auf $30^\circ=150^\circ$ schließen wollte. --- Nehmen wir aber das Symbol $F$ in der weniger allgemeinen Bedeutung, nämlich so, dass es allein die Summe der Reihe $F$ darstellt, so setzen die Überlegungen, durch die wir Gleichung 102 hergeleitet haben, notwendig voraus, dass $y$ nur so weit vom Wert 0 an wächst, bis $x=1$ erreicht ist, also bis $y=\frac12$. Genau an diesem Punkt wird jedoch die \textit{Stetigkeit} der Reihe $P=F(\alpha,\beta,\alpha+\beta+\frac12,4y-4yy)$ unterbrochen, da $\frac{\mathrm dP}{\mathrm dy}$ offenbar plötzlich von einem positiven (endlichen) zu einem negativen Wert springt. \Needspace{12\baselineskip}
In dieser Bedeutung lässt Gleichung 102 also keine Ausdehnung über die Grenzen von $y=\frac12-\BareRoot{\frac12}$ bis $y=\frac12$ hinaus zu. Wenn Sie es vorziehen, können Sie dieselbe Gleichung auch folgendermaßen darstellen:
\[
\begin{aligned}&[103]\quad F(\alpha,\beta,\alpha+\beta+\tfrac12,x)\\&=F(2\alpha,2\beta,\alpha+\beta+\tfrac12,\tfrac{1-\sqrt{1-x}}2)\end{aligned}
\]
oder so:
\[
\begin{aligned}&[104]\quad F(\alpha,\beta,\alpha+\beta+\tfrac12,1-x)\\&=F(2\alpha,2\beta,\alpha+\beta+\tfrac12,\tfrac{1-\BareRoot{x}}2)\end{aligned}
\]
Daraus folgt als Korollar (Formel 48):
\[
\begin{aligned}
{}[105]\quad F(2\alpha,2\beta,\alpha+\beta+\tfrac12,\tfrac12)&=\frac{\Pi(\alpha+\beta-\tfrac12)\Pi(-\tfrac12)}{\Pi(\alpha-\tfrac12)\Pi(\beta-\tfrac12)}\\
&=\frac{\Pi(\alpha+\beta-\tfrac12)\BareRoot{\pi}}{\Pi(\alpha-\tfrac12)\Pi(\beta-\tfrac12)}
\end{aligned}
\]
\SixthArticle{56}
Wendet man Formel 87 auf Gleichung 104 an, so folgt
\[
\begin{aligned}&[106]\quad F(2\alpha,2\beta,\alpha+\beta+\tfrac12,\tfrac{1-\BareRoot{x}}2)\\&=AF(\alpha,\beta,\tfrac12,x)+B\BareRoot{x}\mathbin{.}F(\alpha+\tfrac12,\beta+\tfrac12,\tfrac32,x)\end{aligned}
\]
Daraus ist ersichtlich, dass die Reihe
\[
F(2\alpha,2\beta,\alpha+\beta+\tfrac12,\tfrac{1-t}2)
\]
durch die Reihe
\[
\begin{aligned}&A+Bt+A\frac{\alpha\mathbin{.}\beta}{1\mathbin{.}\tfrac12}\mathbin{.}tt\\&+B\frac{\alpha+\tfrac12\mathbin{.}\beta+\tfrac12}{1\mathbin{.}\tfrac32}t^3\\&+A\frac{\alpha\mathbin{.}\alpha+1\mathbin{.}\beta\mathbin{.}\beta+1}{1\mathbin{.}2\mathbin{.}\tfrac12\mathbin{.}\tfrac32}t^4+\text{usw.}\end{aligned}
\]
dargestellt werden kann, wenn man der Kürze halber Folgendes setzt:
\SourceLeaf{237}
\[
\begin{aligned}&A=\frac{\Pi(\alpha+\beta-\tfrac12)\Pi(-\tfrac12)}{\Pi(\alpha-\tfrac12)\Pi(\beta-\tfrac12)},\\&
B=\frac{\Pi(\alpha+\beta-\tfrac12)\Pi(-\tfrac32)}{\Pi(\alpha-1)\Pi(\beta-1)}\end{aligned}
\]
Hieraus darf man schließen, dass
\[
\begin{aligned}&[107]\quad F(2\alpha,2\beta,\alpha+\beta+\tfrac12,\tfrac{1+\BareRoot{x}}2)\\&=AF(\alpha,\beta,\tfrac12,x)-B\BareRoot{x}\mathbin{.}F(\alpha+\tfrac12,\beta+\tfrac12,\tfrac32,x)\end{aligned}
\]
Sollte jemandem dieser Schluss nicht hinreichend gerechtfertigt erscheinen (obwohl es nicht schwierig wäre, ihn über jeden Zweifel zu erheben), so könnten wir auf folgende Weise zu derselben Gleichung gelangen. Aus Gleichung 87 ergibt sich
\[
\begin{aligned}
&F(2\alpha,2\beta,\alpha+\beta+\tfrac12,\tfrac{1+\BareRoot{x}}2)\\&=CF(2\alpha,2\beta,\alpha+\beta+\tfrac12,\tfrac{1-\BareRoot{x}}2)\\
&\quad+D\left(\frac{1-x}4\right)^{\frac12-\alpha-\beta}F(1-2\alpha,1-2\beta,\tfrac32-\alpha-\beta,\tfrac{1-\BareRoot{x}}2)
\end{aligned}
\]
wenn man der Kürze halber Folgendes setzt:
\[
\begin{aligned}&C=\frac{\Pi(\alpha+\beta-\tfrac12)\Pi(-\tfrac12-\alpha-\beta)}{\Pi(\alpha-\beta-\tfrac12)\Pi(\beta-\alpha-\tfrac12)},\\&
D=\frac{\Pi(\alpha+\beta-\tfrac12)\Pi(\alpha+\beta-\tfrac32)}{\Pi(2\alpha-1)\Pi(2\beta-1)}\end{aligned}
\]
Aus Gleichung 104 lässt sich aber leicht herleiten:
\[
\begin{aligned}
&F(1-2\alpha,1-2\beta,\tfrac32-\alpha-\beta,\tfrac{1-\BareRoot{x}}2)\\&=F(\tfrac12-\alpha,\tfrac12-\beta,\tfrac32-\alpha-\beta,1-x)\\
&=EF(\tfrac12-\alpha,\tfrac12-\beta,\tfrac12,x)\\&\qquad+G\BareRoot{x}\mathbin{.}F(1-\alpha,1-\beta,\tfrac32,x)
\end{aligned}
\]
Dabei setzt man der Kürze halber:
\[
\begin{aligned}&E=\frac{\Pi(\tfrac12-\alpha-\beta)\Pi(-\tfrac12)}{\Pi(-\alpha)\Pi(-\beta)},\\&
G=\frac{\Pi(\tfrac12-\alpha-\beta)\Pi(-\tfrac32)}{\Pi(-\tfrac12-\alpha)\Pi(-\tfrac12-\beta)}\end{aligned}
\]
Hieraus folgt wiederum nach Gleichung 82
\[
\begin{aligned}
&F(1-2\alpha,1-2\beta,\tfrac32-\alpha-\beta,\tfrac{1-\BareRoot{x}}2)\\
&\quad=E(1-x)^{\alpha+\beta-\frac12}F(\alpha,\beta,\tfrac12,x)\\&\qquad+G\BareRoot{x}\mathbin{.}(1-x)^{\alpha+\beta-\frac12}F(\alpha+\tfrac12,\beta+\tfrac12,\tfrac32,x)
\end{aligned}
\]
Setzt man dies ein, so ergibt sich mit
\[
AC+DE2^{2\alpha+2\beta-1}=M,\quad BC+DG2^{2\alpha+2\beta-1}=N
\]
\[
\begin{aligned}&F(2\alpha,2\beta,\alpha+\beta+\tfrac12,\tfrac{1+\BareRoot{x}}2)\\&=MF(\alpha,\beta,\tfrac12,x)+N\BareRoot{x}\mathbin{.}F(\alpha+\tfrac12,\beta+\tfrac12,\tfrac32,x)\end{aligned}
\]
Deren \textit{Form} stimmt mit Gleichung 107 überein. Nun könnten wir zwar allein aus der Natur der Funktion $\Pi$ die Beziehungen $M=A$, $N=-B$ herleiten, da sich nach den Gleichungen 55, 56 leicht zeigen lässt, dass
\[
\begin{aligned}&C=\frac{\cos(\alpha-\beta)\pi}{\cos(\alpha+\beta)\pi},\\&
\frac{DE2^{2\alpha+2\beta-1}}A=-\frac{2\sin\alpha\pi\sin\beta\pi}{\cos(\alpha+\beta)\pi},\\&
\frac{DG2^{2\alpha+2\beta-1}}B=-\frac{2\cos\alpha\pi\cos\beta\pi}{\cos(\alpha+\beta)\pi},\end{aligned}
\]
\SourceLeaf{238}
Doch nicht einmal diese Arbeit ist nötig. Denn setzt man $x=0$, so muss offenbar gelten:
\[
M=F(2\alpha,2\beta,\alpha+\beta+\tfrac12,\tfrac12)=A
\]
Differenziert man dagegen jene Gleichung, so ergibt sich
\[
\begin{aligned}
&x^{-\frac12}\frac{\alpha\beta}{\alpha+\beta+\tfrac12}F(2\alpha+1,2\beta+1,\alpha+\beta+\tfrac32,\tfrac{1+\BareRoot{x}}2)\\
&\quad=2\alpha\beta MF(\alpha+1,\beta+1,\tfrac32,x)\\
&\qquad+\tfrac23(\alpha+\tfrac12)(\beta+\tfrac12)N\BareRoot{x}\mathbin{.}F(\alpha+\tfrac32,\beta+\tfrac32,\tfrac52,x)\\&\qquad+\tfrac12Nx^{-\frac12}F(\alpha+\tfrac12,\beta+\tfrac12,\tfrac32,x)
\end{aligned}
\]
und setzt man darin $x=0$, so folgt
\[
\begin{aligned}
N&=\frac{2\alpha\beta}{\alpha+\beta+\tfrac12}F(2\alpha+1,2\beta+1,\alpha+\beta+\tfrac32,\tfrac12)\\
&=\frac{2\alpha\beta}{\alpha+\beta+\tfrac12}\mathbin{.}\frac{\Pi(\alpha+\beta+\tfrac12)\Pi(-\tfrac12)}{\Pi\alpha\Pi\beta}\\
&=\frac{\Pi(\alpha+\beta-\tfrac12)\Pi(-\tfrac32)}{\Pi(\alpha-1)\Pi(\beta-1)}=-B
\end{aligned}
\]
